\subsection{\bt{Controlled studies and research questions}{受控实验与研究问题}}
\label{sec:metric-setup}
\bi{Together, the official-API comparison and controlled studies ask: does sharing across agents reduce total cost relative to private memory with the same mechanisms (Q1); how the evaluated reuse policy changes exploration and validation work (Q2); how much validation, maintenance, and optimization work can consumers share (Q3); and can this reuse preserve correctness and efficiency after data changes (Q4)? Agents answer with one SQL query. The machine has 48 cores, 98\,GiB RAM, and PostgreSQL 18.6. The population study uses TPC-DS SF1 (24 tables, 2.9 million store-sales rows); held-out metric and data-change studies use synthetic retail data with TPC-DS-style columns (one million store-sales rows, about 396,000 returns and 500,000 catalog sales) and TPC-H SF1.}
{官方API对照与受控实验共同考察：跨智能体共享是否比采用相同机制的私有记忆降低总成本（Q1）；受评估复用策略如何改变探索与验证工作量（Q2）；使用者能够共享多少验证、维护与优化工作（Q3）；这种复用能否在数据变化后保持正确性与效率（Q4）。智能体用一条 SQL 作答。机器为 48 核、98\,GiB 内存，运行 PostgreSQL 18.6。群体实验使用 TPC-DS SF1（24 张表、290 万行门店销售）；留出指标与数据变化实验使用列名仿照 TPC-DS 的合成零售数据（100 万行门店销售、约 39.6 万行退货、50 万行目录销售）及 TPC-H SF1。}

\bi{\emph{Workload.} Five metrics (store revenue, return amount, return rate, catalog revenue, electronics store revenue) each have two learning questions stating the business definition and five scored questions naming only the metric: a new period, a period difference, the highest month, the largest monthly rise, and months above the monthly mean. TPC-DS uses the same seven forms for five template-derived metrics: store sales for one category manager, registered-customer sales, mid-sized-store coupons, returns at known stores, and yearly totals. Reference queries, handwritten or compiled from TPC-DS/TPC-H templates, compute answers on the question's snapshot. Agents and extractor use \llmmodel{} with high reasoning effort, 20 turns, and a 60\,s SQL timeout (120\,s on TPC-DS/TPC-H); methods share tools, with lookup added for memory. Low effort bounds RelBench's repeated feature-search cost; high effort prioritizes exact metric answers, with agent effort fixed across compared arms within each study. In the population study, table descriptions, verified joins, check results, and query results follow the memory scope: per-question without memory, per-agent with private memory, and global with \system. In the held-out metric and data-change studies, memoryless baselines also share these resources, isolating executable definitions and maintenance. Table~\ref{tab:ladder} lists controlled ablations of memory scope and condition maintenance. The condition-sharing microbenchmark reports one run per configuration. Tokens count input plus output; agent time spans the first call to the final answer.}
{\emph{负载。}5 个指标（门店营业额、退货金额、退货率、目录渠道营业额、电子品类门店营业额）各有 2 道给出业务定义的学习题，以及 5 道只给指标名的计分题：新期间、期间差值、最高月、最大月增量、高于月均值的月份数。TPC-DS 的五个模板指标以同样七种形式提问：一位品类经理所管商品的门店销售额、登记顾客销售额、中型门店优惠券金额、关联得上门店的退货金额，以及年度汇总口径。参考查询由人工编写或从 TPC-DS／TPC-H 模板编译，在提问所在快照上计算答案。智能体与提取器使用 \llmmodel{}，推理强度为高，最多 20 轮，SQL 超时 60 秒（TPC-DS／TPC-H 为 120 秒）；各方法共用工具，有记忆的方法增加查找工具。低档位控制 RelBench 反复搜索特征的代价，高档位侧重精确指标作答；各研究内部的比较组保持相同的智能体推理档位。群体实验中，表描述、验证过的连接、检查结果和查询结果遵循记忆作用域：无记忆按题保留，私有记忆按智能体保留，\system\ 全局共享。留出指标与数据变化实验中的无记忆基线也共享这些资源，以隔离可执行定义与维护的收益。表~\ref{tab:ladder} 列出记忆作用域与条件维护的受控消融。条件共享微基准的每种配置报告一次运行。token 为输入与输出之和；智能体耗时从首次调用计到最终答案。}

\begin{table}[t]
% The Times symbol font has no small-caps section sign.
\caption{\bt{Comparison of memory designs and maintenance policies across evaluation studies. P: population study; Q: cost per question; C: cost of sharing; D: data changes.}{各项评估中记忆设计与维护策略的比较。P：群体实验；Q：每道题的代价；C：共享的代价；D：数据变化。}}
\label{tab:ladder}
\footnotesize
\setlength{\tabcolsep}{2.4pt}
\newcommand{\usedin}{$\bullet$}
\begin{tabularx}{\linewidth}{@{}P{0.27\linewidth}Lcccc@{}}
\toprule
\bt{Method}{方法} & \bt{Information available to later sessions}{后续会话可用的信息} & P & Q & C & D\\
\midrule
\bt{No memory}{无记忆} & \bt{nothing; every question is explored}{什么也没有，每道题都探索} & \usedin & \usedin & & \usedin\\
\bt{Single-agent memory}{单智能体记忆} & \bt{\system's memory, but each agent uses only what it learned itself}{\system\ 的记忆，但每个智能体只用自己学到的} & \usedin & & \usedin & \\
\bt{Invalidate on every write}{每次写入即失效} & \bt{definitions, dropped on any write and relearned}{定义，任何写入都丢弃并重新学习} & & & \usedin & \\
\bt{Full recheck per definition}{整定义重查} & \bt{definitions, all conditions rechecked on use, no shared results}{定义，使用时重查全部条件、不共享结果} & & & \usedin & \\
\bt{Validation SQL cache}{验证 SQL 缓存} & \bt{rechecks cached by SQL and dependency versions}{按验证 SQL 与依赖版本缓存重查结果} & & & \usedin & \\
\addlinespace
\bt{\textbf{\system\ (ours)}}{\textbf{\system（本文）}} & \bt{one memory for all agents; definitions checked, maintained, and optimized once}{所有智能体共用一份记忆；定义的检查、维护与优化都只做一次} & \usedin & \usedin & \usedin & \usedin\\
\bottomrule
\end{tabularx}
\end{table}

